% Complete analytic candidate; existing cycles and theorem quantifiers stay.
% No main-manuscript edit and no replacement of a frozen candidate.
\paragraph{Common edges after cycle deletion.}
For cycles $C,D$, put $U=V(C)\cap V(D)\ne\varnothing$. Let $H$
contain their edges on $U$, let $F\subseteq E(H)$ be forbidden, and
let $R\subseteq E(H)\setminus F$ be used by neither cycle.
If the outside vertices induce at most $r_C,r_D$ paths in the
respective cycles, counting isolated vertices and using zero for an
empty set, then
\[
 |(E(C)\cap E(D))\setminus F|
 \ge 2|U|-r_C-r_D-|E(H)|-|F|+|R|.                  \tag{DP}
\]
Deleting $q$ outside paths leaves $|U|-q$ edges on $U$. Removing
forbidden edges loses at most $|F|$ from each cycle; inclusion--exclusion
in $E(H)\setminus(F\cup R)$ gives (DP).

\paragraph{Proof of the safe-common-edge lemma.}
Let $a,b,c,d$ be distinct. Order types $E<M<A$ and reverse coordinates
when necessary. The good targets from $E(aab)$ are $aac,abb,acb,bab,cab$;
the remaining targets from $M(aba)$ are $aca,cba$; those from $A(abc)$
are $abd,adc$. These exhaust changes of one boundary coordinate.
Let $M_0(aba)$ be the image of $D_M$ under
$(0,1,2,3)\mapsto(a,b,c,d)$, and $M_1$ its reversal. Unused-colour
interchange fixes $E,M_0,M_1$, and reversal followed by endpoint-colour
interchange fixes $A$. Thus $E,A$ have one allowed edge set at each
boundary, and $M$ has these two.

For six families use (DP), with $H$ the common recolouring graph and
$F=(P_t\cup P_s)\cap E(H)$. The sufficient upper bounds on
$r_C+r_D$ and lower bounds on $|R|$ are
\[
\begin{array}{c|rrrrr|r}
t\to s&|U|&|E(H)|&|F|&r_C+r_D&|R|&\text{bound}\\\hline
aab\to aac&10&13&2&2&0&3\\
aab\to acb&10&15&1&2&0&2\\
aab\to cab&9&11&1&4&0&2\\
abc\to abd&11&17&0&4&2&3\\
aab\to bab&10&15&0&5&2&2\\
aba\to cba&12&21&0&4&3&2
\end{array}                                                     \tag{CT}
\]
Their fixed-middle layers are respectively
$(C_4,C_4,K_2)$, $(K_2\square K_3,C_4)$, $(K_3,C_4,K_2)$,
$(K_3,C_4,C_4)$, $(K_2\square K_3,K_2,K_2)$, and
$(K_2\square K_3,C_4,K_2)$, linked by four, two, three, six, four,
and seven middle-change edges. These give the vertex and edge counts.
The special edges of $E(aab)$ are $bcd$--$dcd$ and $bdc$--$cdc$:
one from each $E$ side survives in the first common graph, one source
edge in either following graph, and neither in $aab\to bab$.
The deletion identities below give the path bounds.

For $abc\to abd$, both cycles omit $bab$--$cab$ and $bab$--$dab$
from the middle-$a$ triangle: their middle-$a$ paths start at $bab$
by changing the last coordinate to the respective unused colour.
These two edges form $R$.
For $aab\to bab$, the middle-$b$ prism has two row triangles,
with fixed first coordinate $c$ or $d$. The $E$ cycle retains one
edge in either row. In either row, $M_0(bab)$ retains two edges,
including this $E$ edge, while $M_1(bab)$ retains one. The union
therefore uses at most two of the three row edges, leaving at least
one edge unused by both in each row. Choose these two edges for $R$.
For $aba\to cba$, both $M$ cycles have internal degree zero at $bab$
in the middle-$a$ grid. The $A(cba)$ middle-$a$ path starts
$bab,dab$ and ends $bac$, so it also omits $bab$--$bac$ and
$bab$--$bad$. Moreover every middle-change edge of an allowed $M$
or $A$ cycle has the excluded middle-boundary colour $b$ in an
endpoint coordinate. Hence neither cycle uses $dad$--$dcd$.
These three edges form $R$.

For $aab\to abb$, both cycles retain the two edges at $bca$ in the
common middle-$c$ square with opposite corners $bca,dcd$; their
special edges at $dcd$ are absent from the other cycle. For
$abc\to adc$, both retain five edges of the common hexagon
$(bab,bad,cad,dad,dab,cab,bab)$, hence share four. For
$aba\to aca$, the square-support property gives three of four
edges in $Q_d\subset B_{aaa}$ for each cycle, hence two common edges.
The latter two families have no special edges. Every count and support
property above follows from the displayed cycle paths below or
\eqref{eq:path232-cycles}. Colour permutations and reversal preserve
special edges, so the proof applies to every labelled pair of allowed
cycles.

\paragraph{Cycle identities used by the count.}
Brackets record consecutive path vertices. In $E(aab)$, the classes
$w=c,v=c,u=c,u=b$ respectively form
\[
[dbc,cbc,cdc,bdc],\quad[dca,bca,bcd,dcd],\quad
[cbd,cba,cda]\cup[cbc,cdc],\quad[bca,bcd]\cup[bda,bdc].
\]
In $A(abc)$, the classes $v=a,u=b,w=d,u=c$ respectively form
\[
\begin{gathered}
[bab,bad,cad,dad,dab,cab],\quad
[bad,bab,bdb,bcb]\cup[bcd,bca,bda],\\
[bad,cad,dad]\cup[dcd,bcd],\quad[cab,cdb,cda]\cup[cad].
\end{gathered}
\]
In $M_0(aba)$, the classes $u=b,w=b,u=c,w=c$ respectively form
\[
\begin{gathered}
[bcb,bab,bdb]\cup[bad,bcd]\cup[bac,bdc],\quad
[dcb,bcb,bab,bdb,cdb]\cup[cab,dab],\\
[cac,cad,cab]\cup[cdb,cdc],\quad
[cac,bac,bdc,cdc]\cup[dac].
\end{gathered}
\]
For $M_1$, reverse coordinates. Legality lists give every vertex of
each class, and \eqref{eq:path232-cycles} gives exactly these edges.
The middle-$a$ grid of $M_0(aba)$ consists of the path
\[
[bac,cac,cad,cab,dab,dac,dad,bad]
\]
and the internally isolated vertex $bab$; reversal gives $M_1$.
Relabelling this path for $M_i(bab)$ proves the two row-triangle
retention assertions. The same cyclic words show that every inter-grid
edge in $M$ or $A$ has an endpoint coloured $b$. Thus all path and
omission bounds in (CT) hold for every allowed cycle.

% DP itself is general. The path/support/omission properties use the
% original chosen cycles; no arbitrary-Hamilton-cycle claim is made.
% Nine symbolic transition families remain, but the ten-entry common-edge
% certificate and all its target-convention/orbit bookkeeping are removed.
